% Part: many-valued-logic
% Chapter: three-valued-logics
% Section: lukasiewicz

\documentclass[../../../include/open-logic-section]{subfiles}

\begin{document}

\olfileid{mvl}{inf}{god}

\olsection{哥德尔逻辑}

\begin{editorial}
  本节是关于无限值哥德尔逻辑的一个简短占位小节。
\end{editorial}

\begin{defn}\ollabel{def:goedel}
无限值哥德尔逻辑~$\LogGod[\infty]$ 由以下矩阵定义：
\begin{enumerate}
  \item 标准命题语言 $\Lang L_0$，带有联结词
  $\lfalse$、$\lnot$、$\land$、$\lor$、$\lif$。
  \item 真值集为 $V_\infty$。
  \item $1$ 是唯一的指定值，即 $V^+ = \{1\}$。
  \item 真值函数由下列等式给出：
  \begin{align*}
    \tf{\lfalse} & = 0\\
    \tf{\lnot}[\LogGod](x) & = \begin{cases}
      $1$ & \text{若 } x =0\\
      $0$ & \text{否则}
    \end{cases}\\
    \tf{\land}[\LogGod](x,y) & = \min(x,y)\\
    \tf{\lor}[\LogGod](x,y) & = \max(x,y)\\
    \tf{\lif}[\LogGod](x,y) & = \begin{cases}
      1 & \text{若 } x \le y\\
      y & \text{否则。}
    \end{cases}
    \end{align*}
\end{enumerate}
$m$-值哥德尔逻辑的定义方式与此相同，只是 $V = V_m$。
\end{defn}

\begin{prop}
  由 \olref[thr][god]{defn:goedel} 定义的逻辑 $\LogGod[3]$ 与由 \olref{def:goedel} 定义的 $\LogGod[3]$ 是同一个逻辑。
\end{prop}

\begin{proof}
  通过比较 \olref[thr][god]{defn:goedel} 中给出的联结词真值表与 \olref{def:goedel} 中等式所确定的真值表即可看出：
  \begin{center}
    \begin{tabular}{c|c}
      $\tf{\lnot}[\LogGod[3]]$ & \\
      \hline
      $1$ & $0$ \\
      $1/2$ & $0$ \\
      $0$ & $1$
    \end{tabular}
    \quad
    \begin{tabular}{c|ccc}
      $\tf{\land}[\LogGod]$ & $1$ & $1/2$ & $0$ \\
      \hline
      $1$ & $1$ & $1/2$ & $0$ \\
      $1/2$ & $1/2$ & $1/2$ & $0$\\
      $0$ & $0$ & $0$ & $0$
    \end{tabular}
    \\[2ex]
    \begin{tabular}{c|ccc}
      $\tf{\lor}[\LogGod]$ & $1$ & $1/2$ & $0$ \\
      \hline
      $1$ & $1$ & $1$ & $1$ \\
      $1/2$ & $1$ & $1/2$ & $1/2$ \\
      $0$ & $1$ & $1/2$ & $0$
    \end{tabular}
    \quad
    \begin{tabular}{c|ccc}
      $\tf{\lif}[\LogGod]$ & $1$ & $1/2$ & $0$ \\
      \hline
      $1$ & $1$ & $1/2$ & $0$ \\
      $1/2$ & $1$ & $1$ & $0$  \\
      $0$ & $1$ & $1$ & $1$
    \end{tabular}
  \end{center}
\end{proof}

\begin{prop}\ollabel{prop:god-infty-m}
  如果 $\Gamma \Entails[\LogGod[\infty]] !B$，那么对所有 $m \ge 2$，均有 $\Gamma \Entails[\LogGod[m]] !B$。
\end{prop}

\begin{proof}
  练习。
\end{proof}

\begin{prob}
  证明 \olref[mvl][inf][god]{prop:god-infty-m}。
\end{prob}

事实上，其逆命题也成立。

与 $\LogGod[3]$ 一样，$\LogGod[\infty]$ 以所有直觉主义有效的公式作为其重言式，并且同样的非重言式例子在~$\LogGod[\infty]$ 中仍是非重言式：
\begin{align*}
  & p \lor \lnot p && (p \lif q) \lif (\lnot p \lor q) \\
  & \lnot\lnot p \lif p && \lnot(\lnot p \land \lnot q) \lif (p \lor q) \\
  & ((p \lif q) \lif p) \lif p && \lnot(p \lif q) \lif (p \land \lnot q)
\end{align*}
直觉主义无效但却是~$\LogGod[3]$ 的重言式的那个例子，即 $(p \lif q) \lor (q \lif
p)$，在~$\LogGod[\infty]$ 中也是重言式。事实上，
$\LogGod[\infty]$ 可以被刻画为添加了模式 $(!A \lif !B) \lor (!B \lif !A)$ 的直觉主义逻辑。这一结论由 Michael Dummett（杜梅特）证明，因此 $\LogGod[\infty]$ 常被称为哥德尔-杜梅特逻辑~$\Log{LC}$。

\begin{prob}
  证明 $(p \lif q) \lor (q \lif p)$ 是~$\LogGod[\infty]$ 的一个重言式。
\end{prob}

\begin{prob}
  证明 $(p \lif q) \lor (q \lif r) \lor (r \lif s)$ 是 $\LogGod[3]$ 的重言式，但不是~$\LogGod[\infty]$ 的重言式。
\end{prob}

\end{document}